\section{Finite sharpness at width six}

An obstruction from four rows can be an artifact of using too-short
cylinders.  Here that limitation is exact.

\begin{proposition}[Four-row sharpness]\label{prop:four-sharp}
For row order \((A_5,B_7,A_6,B_6)\), the width-five incidence matrix has rank
three and left-kernel vector \((1,1,-1,-1)\).  At width six, the columns
\[
 000021,\quad000023,\quad000210,\quad000231
\]
form the minor
\[
 \begin{pmatrix}
 0&0&1&0\\
 0&1&0&1\\
 1&0&1&0\\
 0&0&0&1
 \end{pmatrix},
 \qquad \det=-1.
\]
Consequently arbitrary totals on these four cycles are realizable by a
width-six cylinder function.
\end{proposition}

The finite sharpness persists after including all earlier obstruction rows.

\begin{proposition}[Cumulative sharpness]\label{prop:cumulative}
For the nine rows
\[
 C_1,A_3,A_4,B_4,A_5,B_5,B_6,A_6,B_7,
\]
there is a selected nine-column width-six incidence minor of determinant
one.  Hence every assignment of nine periodic totals, including their exact
Galois excesses, has a unique interpolation on those selected cylinders.
\end{proposition}

\begin{proof}
The selected columns and full integer matrix are retained in the primary
certificate.  Fraction-free elimination gives determinant one.  An
independent implementation rebuilds the block counters and checks the
four-row determinant without importing the primary code.
\end{proof}

Propositions~\ref{prop:four-sharp} and \ref{prop:cumulative} prevent a common
overclaim.  Theorem~\ref{thm:main} raises the minimum possible local memory
to six; it does not prove that width six works on every primitive orbit.
